\prefacesection{Abstract}

Abstract here (no more than 300 words). An abstract of not more than 300 words must be included and indicates the major emphasis of the thesis, new discoveries, and its contribution to knowledge. The style of abstract varies from discipline to discipline but the student should follow an appropriate style. This page must be numbered ‘iv’.
\\~\\
The Navier-Stokes equations are a set of partial differential equations used to model the behaviour of viscous incompressible fluids. The question of whether or not unique classical solutions of the Navier-Stokes equations, and inviscid Euler equations, exist for all time, given smooth initial conditions, remains unanswered for the three-dimensional problem. 
\\~\\
A singularity is the occurrence of an infinite derivative of the velocity at a single point in the fluid flow. The formation of a singularity from smooth initial conditions would invalidate the equations as accurate descriptions of fluid mechanics. The Euler-Voigt equations are an inviscid regularization of the Euler equations which are known to not form singularities from smooth initial conditions. However, by decreasing the regularization parameter we approximate Euler flows. We examine the Euler-Voigt equations with a PDE-constrained optimization problem to identify initial conditions that produce 'extreme' behaviour in numerical simulations.  \#\# We present numerical evidence that CONCLUSION \#\#